% ============================================================================
% Speichermanagement: Einverstaendniserklaerung und Information zur
% Datennutzung fuer neue Mitglieder (eine Seite, leicht verstaendlich).
% Bauen: pdflatex speichermanagement-einverstaendnis.tex
%
% Standardweg ist die Online-Zustimmung im Mitgliederbereich
% (/user/<nr>/speichermanagement, Komponente ConsentText.svelte); dieses
% Blatt ist die inhaltsgleiche Papier-Variante (Ausdruck-Fallback). Bei
% Textaenderungen BEIDE Stellen anpassen und die Version in
% website/src/lib/consent/speichermanagement.js erhoehen.
% Fakten: Batteriemanagement/openhab/setup/README.md (Abschnitt Status-Push),
% Datenschutzerklaerung website/src/routes/(website)/datenschutz/
% ============================================================================
\documentclass[a4paper,10pt]{article}
\usepackage[utf8]{inputenc}
\usepackage[T1]{fontenc}
\usepackage[ngerman]{babel}
\usepackage[margin=1.9cm]{geometry}
\usepackage{lmodern}
\usepackage{parskip}
\usepackage{enumitem}
\usepackage{xcolor}
\usepackage[colorlinks=true,linkcolor=black,urlcolor=blue!50!black]{hyperref}
\pagestyle{empty}

\begin{document}

\begin{center}
    {\LARGE\bfseries Speichermanagement}\\[0.3em]
    {\large Information zur Datennutzung und Einverständniserklärung}\\[0.3em]
    {\small Verantwortlicher: Erneuerbare-Energie-Gemeinschaft ISCHLSTROM,
     ZVR-Zahl 1915125185, Faschlweg 39, 4820 Bad Ischl,
     info@ischlstrom.org \(\cdot\) Stand: August 2026}
\end{center}

\section*{Was das Speichermanagement macht}

Sie bekommen von uns einen kleinen Computer (Raspberry Pi), der bei Ihnen
neben dem Router steht. Er sorgt dafür, dass Ihr Batteriespeicher bevorzugt
dann lädt, wenn in der Energiegemeinschaft Sonnenstrom übrig ist. So bleibt
mehr Strom in der Gemeinschaft und Sie sparen Netzstrom. Die
Sicherheitsfunktionen Ihres Wechselrichters bleiben dabei unverändert, und
Sie können das Gerät jederzeit einfach abstecken.

\section*{Welche Daten übertragen werden}

Damit die Steuerung funktioniert und wir Ihnen bei Problemen helfen können,
meldet das Gerät laufend an unseren Server:

\begin{itemize}[nosep]
    \item den \textbf{Ladestand} Ihres Batteriespeichers,
    \item die aktuelle \textbf{Leistung} von Photovoltaik, Netz und Batterie,
    \item den \textbf{technischen Zustand} des Geräts selbst
          (z.\,B. Temperatur, anstehende Updates, Fehlermeldungen).
\end{itemize}

\textbf{Nicht übertragen} werden Verbrauchsdaten Ihres Haushalts oder
Informationen über einzelne Geräte in Ihrem Haus. Falls Ihr
Wechselrichter ein Passwort braucht (Fronius GEN24), wird es
verschlüsselt gespeichert, einmalig an Ihr Gerät übertragen und danach
am Server gelöscht.

\section*{Wer die Daten sieht und wie lange wir sie speichern}

Ihre Daten sehen nur Sie selbst (im Mitgliederbereich auf ischlstrom.org)
und der Vorstand (zur Betreuung der Anlagen). Alles liegt auf den eigenen
Servern des Vereins in der EU; wir geben nichts an Dritte weiter.
Verlaufsdaten löschen wir automatisch nach 30 Tagen. Für die Wartung hat
der Vorstand einen gesicherten Fernzugriff auf das Gerät; er wird nur
genutzt, um die Anlage in Betrieb zu halten.

\section*{Freiwilligkeit und Widerruf}

Die Teilnahme ist freiwillig. Rechtsgrundlage ist Ihre Einwilligung
(Art. 6 Abs. 1 lit. a DSGVO). Sie können sie jederzeit ohne Angabe von
Gründen widerrufen, formlos per E-Mail an
\href{mailto:info@ischlstrom.org}{info@ischlstrom.org} oder auf Ihrer
Seite \emph{Speichermanagement} auf ischlstrom.org. Wir stellen die
Steuerung und die Datenübertragung dann ein und löschen Ihre gespeicherten
Daten. Der Widerruf berührt nicht die Rechtmäßigkeit der bis dahin
erfolgten Verarbeitung. Ihnen stehen außerdem die Rechte auf Auskunft,
Berichtigung, Löschung, Einschränkung der Verarbeitung und
Datenübertragbarkeit zu; Beschwerden richten Sie an uns oder an die
Österreichische Datenschutzbehörde (\url{https://www.dsb.gv.at}). Alle
Einzelheiten finden Sie in unserer Datenschutzerklärung:
\url{https://ischlstrom.org/datenschutz}

\section*{Einverständniserklärung}

Ich habe die Informationen oben gelesen und bin damit einverstanden, dass
die ISCHLSTROM Energiegemeinschaft meinen Batteriespeicher wie beschrieben
steuert und die genannten Daten dafür verarbeitet.

\vspace{2em}
\noindent
\begin{tabular}{@{}p{0.45\textwidth}p{0.45\textwidth}@{}}
    \rule{0.4\textwidth}{0.4pt} & \rule{0.4\textwidth}{0.4pt}\\
    Name (in Blockschrift) & Mitgliedsnummer\\[2em]
    \rule{0.4\textwidth}{0.4pt} & \rule{0.4\textwidth}{0.4pt}\\
    Ort, Datum & Unterschrift\\
\end{tabular}

\end{document}
